% Part A (tech core) of the AI-for-X seminars: prompts and the CRAFT framework. 8 frames.
% Counterpart of ai_core_nontech_prompts (same size, same CRAFT frames, engineering examples).
% Part A is shared: edit here, then recompile every deck that includes it. Trim by commenting frames, not deleting.

%%%%%%%%%%%%%%%%%%%%%%%%%%%%%%%%%%%%%%%%%%%%%%%%%%%%%%%%%%
\begin{frame}[fragile]\frametitle{What Is a Prompt?}
\begin{itemize}
\item A \textbf{prompt} is the input text that conditions the model's next tokens. Everything in the context window counts as prompt.
\item Chat APIs split it into messages: a \textbf{system} message (standing instructions), \textbf{user} messages, and the model's earlier replies.
\item More specific text narrows the distribution of likely answers, so the result is more specific.
\item \textbf{Prompt engineering} is designing, testing and refining that input until the output is reliably usable.
\end{itemize}
\end{frame}

%%%%%%%%%%%%%%%%%%%%%%%%%%%%%%%%%%%%%%%%%%%%%%%%%%%%%%%%%%
\begin{frame}[fragile]\frametitle{Same Task, Different Prompts}
\begin{itemize}
\item Same model, same task: a vague prompt gives a generic answer, a specific one gives a usable answer.
\item A chat keeps the earlier messages in context, so reply to refine instead of starting again.
\end{itemize}
\begin{lstlisting}[language={}]
Vague:      Write code to parse the log file.

Specific:   Write a Python function that reads a web
            server log, and returns the count of 5xx
            responses per hour, as a dict.

Follow-up:  Add type hints and a test with 3 sample lines.
\end{lstlisting}
\end{frame}

%%%%%%%%%%%%%%%%%%%%%%%%%%%%%%%%%%%%%%%%%%%%%%%%%%%%%%%%%%
\begin{frame}[fragile]\frametitle{A Framework for Every Prompt: CRAFT}

\begin{center}
\Large \textbf{C R A F T}
\end{center}

\medskip
\begin{itemize}
    \item \textbf{C}ontext, who you are, what the situation is
    \item \textbf{R}ole, who the AI should act as
    \item \textbf{A}ction, exactly what you want it to do
    \item \textbf{F}ewshots, one or two examples of what ``good'' looks like
    \item \textbf{T}weaks, tone, length, format, and what to avoid
\end{itemize}

\medskip
\begin{center}
\textbf{Run through the five letters in your head before you type. That is the whole framework.}
\end{center}

\end{frame}

%%%%%%%%%%%%%%%%%%%%%%%%%%%%%%%%%%%%%%%%%%%%%%%%%%%%%%%%%%
\begin{frame}[fragile]\frametitle{CRAFT: Context and Role}

\begin{itemize}
    \item \textbf{Context}, without it, the model is guessing: the language and version, the data, the constraint, what has already been tried.
    \item \textbf{Role}, a role steers the register and the checks applied: a ``security reviewer'' reads code differently from a ``performance engineer.''
\end{itemize}

\begin{lstlisting}[language=HTML]
Weak:  Review this function.

CRAFT: You are a security reviewer (Role).
       The function handles user uploads in a Flask
       app on Python 3.11 (Context).
\end{lstlisting}

\end{frame}

%%%%%%%%%%%%%%%%%%%%%%%%%%%%%%%%%%%%%%%%%%%%%%%%%%%%%%%%%%
\begin{frame}[fragile]\frametitle{CRAFT: Action and Fewshots}

\begin{itemize}
    \item \textbf{Action}, one clear instruction. ``Look at this'' is not an action, ``List up to 5 issues, most serious first'' is.
    \item \textbf{Fewshots}, show, do not just tell. One example of the output you want often does more than a paragraph of description, and it fixes the format.
\end{itemize}

\begin{lstlisting}[language=HTML]
Action:   List up to 5 security issues, most serious
          first.

Fewshot:  Match this format:
          {"line": 12, "issue": "SQL built from input",
           "fix": "use a parameterised query"}
\end{lstlisting}

\end{frame}

%%%%%%%%%%%%%%%%%%%%%%%%%%%%%%%%%%%%%%%%%%%%%%%%%%%%%%%%%%
\begin{frame}[fragile]\frametitle{CRAFT: Tweaks}

\begin{itemize}
    \item \textbf{Tweaks} are the small dials that turn a generic answer into a usable one:
    \begin{itemize}
        \item \textbf{Tone}, formal / terse / explanatory
        \item \textbf{Length}, word count or ``under 10 lines''
        \item \textbf{Format}, bullet list / table / JSON / code only
        \item \textbf{Constraints}, what to avoid, such as ``no new dependencies''
    \end{itemize}
    \item Add tweaks last. When you call a model by API, temperature is one more dial: lower for repeatable output, higher for variety.
\end{itemize}

\begin{lstlisting}[language=HTML]
Tweaks: Output only valid JSON, no prose, no new
        dependencies.
\end{lstlisting}

\end{frame}

%%%%%%%%%%%%%%%%%%%%%%%%%%%%%%%%%%%%%%%%%%%%%%%%%%%%%%%%%%
\begin{frame}[fragile]\frametitle{CRAFT, All Together}

\begin{lstlisting}[language=HTML]
You are a security reviewer (R). The function below
handles user uploads in a Flask app on Python 3.11 (C).

List up to 5 security issues, most serious first (A).

Match this format:
{"line": 12, "issue": "SQL built from input",
 "fix": "use a parameterised query"}              (F)

Output only valid JSON, no prose (T).

[FUNCTION CODE]
\end{lstlisting}

\end{frame}

%%%%%%%%%%%%%%%%%%%%%%%%%%%%%%%%%%%%%%%%%%%%%%%%%%%%%%%%%%
\begin{frame}[fragile]\frametitle{CRAFT: The Takeaways}

\begin{itemize}
    \item Not every prompt needs all five letters spelled out explicitly, but running through them in your head catches what you would otherwise forget.
    \item Treat a prompt like code: keep the ones that work, version them, and test them on several inputs, not one.
    \item \textbf{Always remember:} review the output before using it, and never paste secrets or confidential data into a public AI tool.
\end{itemize}

\end{frame}
